


%%%%%%%%%%%%%%%
\begin{frame}{Pour n'a-t-on pas perdu d'information?}
 
Parce qu'on peut reconstituer \alert{exactement} la relation complète
\vfill

    
\begin{redbullet}
  \item Connaissant \texttt{id\_expert}, je connais 
     toutes les autres informations sur l'expert.
 \item Connaissant \texttt{id\_manuscrit}, je connais \texttt{id\_expert}
     \item  Par \alert{transitivité}, connaissant \texttt{id\_manuscrit}, je  connais \texttt{id\_expert},
     et connaissant \texttt{id\_expert}, je connais toutes les autres informations sur l'expert
\end{redbullet}

\vfill


\begin{blocgauche}
L'association de la clé primaire et de la clé étrangère est le mécanisme
de liaison des nuplets en relationnel.
\end{blocgauche}

\end{frame}


%%%%%%%%%%%%%%%
\begin{frame}{Avec un exemple}
 
\begin{tabular}{c|c|c}
\textbf{id\_expert} & \textbf{nom} & \textbf{adresse}\\ \hline
   1 &Sophie &rue Montorgueil, Paris\\
   2& Philippe  &rue des Martyrs, Paris\\
\hline
\end{tabular}

\vfill

\begin{tabular}{c|l|p{4cm}|c|p{3cm}}
\textbf{id\_manuscrit} & \textbf{nom} & \textbf{titre} & \textbf{id\_expert}  & \textbf{Commentaire}\\ \hline
   10& Serge&  L'arpète &2& Une réussite\\
   20& Cécile& Un art du chant grégorien sous Louis XIV& 1& Bravo\\
\hline
\end{tabular}

\vfill

\begin{blocgauche}
La valeur d'une clé étrangère est \alert{toujours} celle
d'une clé de la table référencée.
\end{blocgauche}
\end{frame}


%%%%%%%%%%%%%%%
\begin{frame}{Schéma normalisé et  contraintes}
 
On normalise les schémas, et on demande au système de garantir deux contraintes.

\begin{defblock}{Contrainte d'unicité}
Une valeur de clé ne peut apparaître \alert{qu'une fois} dans une relation.
\end{defblock}
 
\vfill

\begin{defblock}{Contrainte d'intégrité référentielle}
La valeur d'une clé étrangère doit \alert{toujours}
   être également une des valeurs de la clé référencée.\end{defblock}

\vfill

\begin{blocgauche}
On obtient des schémas sans anomalie.
\end{blocgauche}

\end{frame}

%%%%%%ù

\begin{frame}{À retenir}

Un schéma de BD correct:
\vfill
\begin{redbullet}
  \item Toute relation a une clé
 \item Tous les attributs dépendent \alert{directement} de la clé (forme normale)
 \item On peut lier les nuplets les uns aux autres avec des clés étrangères.
 \item Le système évite les anomalies en garantissant
    la \alert{contrainte d'unicité} (clé) et la \alert{contrainte d'intégrité référentielle} (clé étrangère)
\end{redbullet}

\vfill

\alert{Comment trouve-t-on le bon schéma\,?} Nous y reviendrons.
\end{frame}

